\subsection{Important Details of Convolution}\label{sec:important-details-of-convolution}

\subsubsection{Correlation vs. Mathematical Convolution}\label{sec:correlation-vs-mathematical-convolution}

This section is mainly a \textbf{terminology clarification}. Up to now, we have been calling the sliding filter operation a \textbf{convolution}, but mathematically the operation written in the earlier sections is actually a \textbf{correlation}. The two operations are \textbf{almost identical}. The only difference is whether the filter is \textbf{used \emph{as it is}} or \textbf{\emph{flipped}} before being applied.

\highspace
For \textbf{correlation}, we write the operation as (\autopageref{eq:linear-filter-rgb}):
\begin{equation*}
    T[I](r, c) = \sum_{u,v \in U} w(u, v) \cdot I(r + u, c + v)
\end{equation*}
Here, the filter weights are used directly: $w(u, v)$.

\highspace
For mathematical \textbf{convolution}, instead, we flip the filter before applying it:
\begin{equation*}
    T[I](r, c) = \sum_{u,v \in U} w(-u, -v) \cdot I(r + u, c + v)
\end{equation*}
Here, the filter weights are flipped: $w(-u, -v)$. Or equivalently, we can flip the image instead of the filter:
\begin{equation*}
    T[I](r, c) = \sum_{u,v \in U} w(u, v) \cdot I(r - u, c - v)
\end{equation*}
So the image indexing can be written in exactly the same form as correlation, provided that we replace $w(u, v) \to w(-u, -v)$. Thus, the \hl{convolution correspond to using the \textbf{flipped filter}.}

\begin{figure}[!htp]
    \centering

    \definecolor{kernelbg}{RGB}{242, 245, 250}
    \definecolor{flipbg}{RGB}{253, 238, 238}
    \definecolor{cblue}{RGB}{31, 110, 165}
    \definecolor{cred}{RGB}{180, 40, 30}
    \definecolor{corange}{RGB}{210, 80, 0}
    \definecolor{cteal}{RGB}{15, 140, 115}
    \definecolor{darknavy}{RGB}{20, 35, 60}

    \begin{tikzpicture}[
        >=Stealth,
        cell/.style={
            draw=black!60,
            rectangle,
            minimum size=0.95cm,
            anchor=center,
            inner sep=0pt,
            font=\small
        }
    ]

    % =============================================================
    % MATRIX 1: ORIGINAL KERNEL w(u,v)
    % =============================================================
    \matrix (K1) [
        matrix of math nodes, 
        nodes={cell, fill=kernelbg},
        column sep=-\pgflinewidth,
        row sep=-\pgflinewidth
    ] {
        |[fill=cblue!35, text=cblue, font=\bfseries]| 1 & |[fill=corange!20]| 2 & |[fill=corange!35, text=corange, font=\bfseries]| 3 \\
        4 & 5 & 6 \\
        |[fill=cteal!35, text=cteal, font=\bfseries]| 7 & |[fill=cteal!20]| 8 & |[fill=cred!35, text=cred, font=\bfseries]| 9 \\
    };

    % Coordinate Axes for K1 (Pure black for crisp legibility)
    \node[above=2pt of K1-1-1.north, font=\tiny, text=black] {$v=-1$};
    \node[above=2pt of K1-1-2.north, font=\tiny, text=black] {$v=0$};
    \node[above=2pt of K1-1-3.north, font=\tiny, text=black] {$v=1$};

    \node[left=3pt of K1-1-1.west, font=\tiny, text=black] {$u=-1$};
    \node[left=3pt of K1-2-1.west, font=\tiny, text=black] {$u=0$};
    \node[left=3pt of K1-3-1.west, font=\tiny, text=black] {$u=1$};

    % =============================================================
    % MATRIX 2: FLIPPED KERNEL w(-u,-v)
    % =============================================================
    \matrix (K2) [
        right=3.5cm of K1, 
        matrix of math nodes, 
        nodes={cell, fill=flipbg},
        column sep=-\pgflinewidth,
        row sep=-\pgflinewidth
    ] {
        |[fill=cred!35, text=cred, font=\bfseries]| 9 & |[fill=cteal!20]| 8 & |[fill=cteal!35, text=cteal, font=\bfseries]| 7 \\
        6 & 5 & 4 \\
        |[fill=corange!35, text=corange, font=\bfseries]| 3 & |[fill=corange!20]| 2 & |[fill=cblue!35, text=cblue, font=\bfseries]| 1 \\
    };

    % Coordinate Axes for K2
    \node[above=2pt of K2-1-1.north, font=\tiny, text=black] {$-v=1$};
    \node[above=2pt of K2-1-2.north, font=\tiny, text=black] {$-v=0$};
    \node[above=2pt of K2-1-3.north, font=\tiny, text=black] {$-v=-1$};

    \node[right=3pt of K2-1-3.east, font=\tiny, text=black] {$-u=1$};
    \node[right=3pt of K2-2-3.east, font=\tiny, text=black] {$-u=0$};
    \node[right=3pt of K2-3-3.east, font=\tiny, text=black] {$-u=-1$};

    % =============================================================
    % MATRIX TITLES
    % =============================================================
    \node[above=0.8cm of K1, font=\bfseries\normalsize, text=darknavy] (T1) {Original Kernel $w(u,v)$};
    \node[above=0.8cm of K2, font=\bfseries\normalsize, text=darknavy] (T2) {Flipped Kernel $w(-u,-v)$};

    % =============================================================
    % CENTRAL TRANSFORMATION BRIDGE
    % =============================================================

    % Curved Rotation Arc
    \draw[->, line width=1.2pt, cblue, shorten >=2pt, shorten <=2pt] 
        ($(K1.east)+(0.25cm, 0.65cm)$) to[out=30, in=150] 
        node[midway, above=3pt, font=\bfseries\small, text=darknavy] {$180^\circ$ Spatial Rotation} 
        ($(K2.west)+(-0.25cm, 0.65cm)$);

    % Middle Row Mapping Arrow
    \draw[->, line width=1.1pt, cblue!90!black, shorten <=6pt, shorten >=6pt] 
        (K1-2-3.east) -- (K2-2-1.west)
        node[midway, fill=white, inner sep=2pt, font=\scriptsize\itshape, text=black] {$w(u,v) \mapsto w(-u,-v)$};

    % =============================================================
    % INDEX MAPPING ANNOTATIONS BELOW
    % =============================================================
    \node[below=0.45cm of K1, font=\scriptsize, text=black, align=center] {
        Top-Left: $w(-1,-1) = 1$\\[.3em]
        Bottom-Right: $w(1,1) = 9$
    };

    \node[below=0.45cm of K2, font=\scriptsize, text=black, align=center] {
        Top-Left: $w_{\text{flip}}(-1,-1) = w(1,1) = 9$\\[.3em]
        Bottom-Right: $w_{\text{flip}}(1,1) = w(-1,-1) = 1$
    };

    \end{tikzpicture}

    \caption{Kernel flipping in mathematical convolution. Replacing $w(u,v)$ with $w(-u,-v)$ reverses the kernel along both spatial dimensions, which is equivalent to a $180^\circ$ rotation. Thus, while correlation applies $w(u,v)$ directly, mathematical convolution applies its spatially flipped version $w(-u,-v)$.}
    \label{fig:correlation-vs-convolution}
\end{figure}

\highspace
\begin{flushleft}
    \textcolor{Green3}{\faIcon{question-circle} \textbf{Why this distinction?}}
\end{flushleft}
This distinction is important in different contexts:
\begin{itemize}
    \item \important{Classical Signal/Image Processing, the distinction matters}. Suppose we manually choose a fixed kernel:
    \begin{equation*}
        w = \begin{bmatrix}
            1 & 0 & -1 \\
            1 & 0 & -1 \\
            1 & 0 & -1
        \end{bmatrix}
    \end{equation*}
    Correlation uses exactly this matrix. True convolution first flips it:
    \begin{equation*}
        w_{\text{flip}} = \begin{bmatrix}
            -1 & 0 & 1 \\
            -1 & 0 & 1 \\
            -1 & 0 & 1
        \end{bmatrix}
    \end{equation*}
    Those can give different responses. For example, opposite signs for an oriented edge. So if the kernel is \textbf{given beforehand}, we absolutely need to know which operation we are performing.

    \highspace
    Mathematical convolution specifically includes the flip because that definition gives convolution its \hl{useful algebraic structure in signal processing, particularly through the Fourier theorem.}


    \item \important{In a CNN, however, nobody gives us $w$}. Suppose true convolution would need to learn:
    \begin{equation*}
        w = \begin{bmatrix}
            1 & 0 & -1 \\
            1 & 0 & -1 \\
            1 & 0 & -1
        \end{bmatrix}
    \end{equation*}
    If our library performs correlation instead, training can simply learn the flipped version:
    \begin{equation*}
        w_{\text{flip}} = \begin{bmatrix}
            -1 & 0 & 1 \\
            -1 & 0 & 1 \\
            -1 & 0 & 1
        \end{bmatrix}
    \end{equation*}
    Then $\text{corr}(x, w_{\text{flip}})$ \textbf{produces the same result} that $\text{conv}(x, w)$ would have produced. So from the optimizer's perspective, \hl{learn $w$ and flip it during convolution}, and \hl{learn the already-flipped $w_{\text{flip}}$ and use correlation}, are \textbf{equivalent} parameterizations.

    \highspace
    That is what Goodfellow et al.\cite{Goodfellow-et-al-2016} means when saying that an algorithm using kernel flipping simply learns a kernel that is flipped relative to the one learned by an algorithm without flipping.

    \begin{deepeningbox}[: What is Goodfellow's message?]
        Goodfellow means that \hl{two CNNs can learn the same operation even if one implements true convolution and the other implements correlation; their stored kernels will simply be flipped versions of each other.}

        \highspace
        Suppose the useful operation requires, under \textbf{true convolution}, the learned kernel:
        \begin{equation*}
            W = \begin{bmatrix}
                1 & 2 & 3 \\
                4 & 5 & 6 \\
                7 & 8 & 9
            \end{bmatrix}
        \end{equation*}
        True convolution flips it before applying it:
        \begin{equation*}
            W' = \begin{bmatrix}
                9 & 8 & 7 \\
                6 & 5 & 4 \\
                3 & 2 & 1
            \end{bmatrix}
        \end{equation*}
        So the input actually gets multiplied against the pattern $W'$.

        \highspace
        Now consider another CNN implementation that performs \textbf{correlation}. It doesn't flip anything, so during training, it can simply learn the flipped version of the kernel:
        \begin{equation*}
            W' = \begin{bmatrix}
                9 & 8 & 7 \\
                6 & 5 & 4 \\
                3 & 2 & 1
            \end{bmatrix}
        \end{equation*}
        Then it applies $W'$ directly. Therefore:
        \begin{equation*}
            \underbrace{\text{Conv}(X, W)}_{\text{flip } W, \text{ then apply}} = \underbrace{\text{Corr}(X, W')}_{\text{apply } W' \text{ directly}}
        \end{equation*}
        In other words:
        \begin{equation*}
            W_{\text{convolutional implementation}} = \text{flip}\left(W_{\text{correlation implementation}}\right)
        \end{equation*}
        The \textbf{numerical values stored as the kernel are different}, but the operation performed on the input can be identical. This is possible precisely because \textbf{the kernel is not predetermined}. Gradient descent is free to learn whichever orientation is needed.
        \begin{equation*}
            \begin{array}{c}
                \textbf{True convolution}\\[2mm]\text{learn }W\\\downarrow\\\text{flip }W\\\downarrow\\\text{apply}
            \end{array}\qquad\equiv\qquad
            \begin{array}{c}
                \textbf{CNN correlation}\\[2mm]\text{learn }\operatorname{flip}(W)\\\downarrow\\\text{no flip}\\\downarrow\\\text{apply}
            \end{array}
        \end{equation*}
        So Goodfellow is \textbf{not saying that training explicitly takes a learned kernel and flips it}. Rather, \hl{if we trained two otherwise equivalent systems, one using mathematical convolution and one using correlation, the kernels they learn would correspond to each other through a flip. That's why the choice does not reduce what the CNN can learn.}
    \end{deepeningbox}
\end{itemize}

\newpage

\begin{takeawaysbox}
    \begin{itemize}
        \item Correlation applies the filter directly, whereas mathematical convolution flips the filter before applying it.
        \item CNN libraries typically implement correlation but conventionally call it convolution.
        \item Since CNN filters are learned from data, this distinction is usually irrelevant in practice, provided that the same convention is used consistently during learning and inference.
        \item Correlation: $I(r + u, c + v)$.
        \item Convolution: $I(r - u, c - v)$.
    \end{itemize}
\end{takeawaysbox}